\section{Active alternate proof attacks and resolved bridges}

The central construction, intrinsic nearby-cycle/Leray proof, based van
Kampen calculation, sphere recognition, and algebraic-dimension calculation
are now verified in this workbench.  The following list records the remaining
alternate-collapse tests, the resolved based-map test, and one abstract
conditional Section~10 conductor theorem; its application to the reconstructed
candidate is verified.  Matching ranks is insufficient.

\subsection{Cusp quotient and collapse}

The source defines an infinite toric threefold $Y$ from the height-one cones over the periodic $A_2$ triangulation and an action of the form
\[
\Psi_{\bar\lambda}(x,t)
=\bigl(c_{\bar\lambda}(t)t^{B_0\bar\lambda}x,t\bigr),
\qquad \bar\lambda\in\bar\Lambda.
\]

\begin{obligation}[C1: finite equivariant collapse model]
\Open\ The locally finite fan, smooth toric charts, action law, bounded
chart estimates, proper quotient, and fibre geometry are verified in
\cref{thm:candidate-cusp-geometry}.  Audit the separate
topological claims in Proposition~7.2 and Lemmas~7.5--7.10: the finite
quotient model and the equivariant diagram
\[
\begin{CD}
Y_{\le\eta} @>{r_Y}>> Y_0\\
@VV{/\bar\Lambda}V @VV{/\bar\Lambda}V\\
N'_0 @>{r}>> W.
\end{CD}
\]
The strong deformation retraction and its continuity across every stratum
remain an independent certificate, not a premise of the Leray proof.
\end{obligation}

\begin{obligation}[C2: integral collapse matrix]
\Open\ Audit Proposition~7.12, its use of Lemma~A.2(6), and the normalizations
left outside Appendix~B by Remark~B.4.  With fixed cellular orientations,
validate or refute the exact sequence
\[
0\longrightarrow
\operatorname{im}(\wedge^qM_0-I)
\longrightarrow\wedge^q\Lambda
\xrightarrow{c_*}H_q(W;\mathbb Z)
\longrightarrow0.
\]
In particular, derive $c_*\colon\Lambda\to H_1(W)=\bar\Lambda$ as the quotient map, prove
\[
c_*(\widehat\gamma\wedge\widehat u)=[\mathbb T_B],
\qquad
c_*(\widehat w\wedge\widehat\delta)=0,
\]
and record the degree-three orientations rather than inferring them from ranks.
\end{obligation}

\begin{obligation}[C3: collapse equals specialization]
\Open\ Theorem~B.1 and the integral saturation argument are verified in
\cref{thm:p0-integral-specialization}.  Compare that intrinsic map with
Proposition~7.12 and prove or refute, functorially and integrally, that
\[
c^*=i_F^*r^*=\operatorname{sp}_{\mathrm{nc}}^q,
\qquad
\operatorname{im}c^*=(\wedge^qV)^{T_0},
\]
including the degree-three sign.  Saturation and torsion of specialization
are already fixed; the comparison diagram must state the precise nearby-cycle
edge map, and equality may not be defined by choosing the manuscript's answer
as a normalization.
\end{obligation}

\subsection{Based inclusions}

Choose one basepoint on the zero section and fixed whiskers through every collar.

\begin{obligation}[V1: primitive meridians and drift]
\Verified\ The actual coordinate gluings, one fixed zero-section base locus,
and explicit whiskers have no unrecorded conjugations and give
\[
\widehat g_j^{m_j}=t_{v_j},
\qquad
\widetilde a_0=t_\mu\widehat\rho_2^{-1}\widehat\rho_1^{-1}.
\]
The proof verifies that the clockwise meridian lifts to
$\widehat g_j$ rather than its inverse or a nonprimitive power, that the two
signs agree, that the zero-section whiskers add no hidden translation, and
that the tautological cusp disc gives $\mu=0$; see
\cref{prop:geometric-pi1}.  V1 is retained here as a resolved audit record.
\end{obligation}

\subsection{Mayer--Vietoris and Leray}

\begin{obligation}[MV1: geometric boundary cylinders]
\Open\ Audit Lemmas~7.18--7.21 and Theorem~7.22.  For
$K_k=\ker(H_k(M)\to H_k(W))$, derive geometrically
\[
K_1=\mathbb Z[\widetilde a_0],
\qquad
K_2=\langle\xi_{\widehat w},\xi_{\widehat\delta}\rangle,
\]
and fix one orientation convention for
\[
\partial(j_*\xi_\nu)=(I-A_1^{-1})\nu=(I-A_2)\nu.
\]
The same calculation must recover the multiple-fibre class
\([F]\longmapsto(3[S_1],-4[S_2])\) and the index-two degree-three image at $p_2$.
\end{obligation}

\begin{obligation}[L1: remaining transgression comparison]
\Open\ The Leray transgressions themselves are verified in
\cref{lem:first-transgression,lem:higher-transgressions}, including their
identification with the Cech discrepancy through the canonical truncation
triangle.  What remains is only the comparison with the separately oriented
Mayer--Vietoris obstruction.  For reference, the proved discrepancy triple
for $12\gamma$ is
\[
\bigl(4\ell_1-\alpha_1,\ 3\ell_2-\alpha_2,\
-(12\ell_0-\alpha_1-\alpha_2)\bigr),
\]
whose sum is $-p$ and determines the sign in
\[
d_2^{0,1}(12\gamma)=\pm p\,\omega.
\]
The cup-product and coinvariant identifications are already proved natural in
the Leray route; the remaining task is compatibility with the separately
oriented Mayer--Vietoris matrix.
\end{obligation}

\subsection{Uniform non-torsion application}

\begin{obligation}[P10.8: one line bundle non-torsion on every fibre]
\Verified\ The candidate's period/countability route now gives a single
\(L\in\operatorname{Pic}(X)\) whose restriction is non-torsion on every
scheme-theoretic fibre and on both reduced finite fibres.  The workbench
proves
\[
R^1f_*\mathcal O_X\simeq\mathcal O\oplus\mathcal O(-1),
\qquad \operatorname{Pic}(X)\simeq\mathbb C,
\]
degree-one base change, and the exact orders
\(\operatorname{ord}(\mathcal O_X(S_j)|_{S_j})=m_j\).  The equation
\(c\tau+d\mu-b=0\) then yields the simultaneous countability conclusion.
The cusp extension is reconstructed from intrinsic nearby cycles and does
not use C1--C3 or MV1.  See
\cref{thm:structure-direct-images-audit,thm:uniform-nontorsion-audit}.
\end{obligation}

The intrinsic nearby-cycle/Leray route proves the recognition statement
\cref{cor:candidate-sphere}.  C1--C3, MV1, and the comparison part of L1
remain active alternate tests; V1 is resolved.  C1--C3 and MV1 test the
manuscript's stronger claims about the explicit collapse map, not a missing
hypothesis of the verified sphere proof.
